\subsection{The same actual normalized-model line in its original polarization}

\begin{lemma}[The original product inclusion retains the whole original projection]
\label{local:actual-product-polynomial-original-inclusion-square}
\lean{EnvelopeProperty.ProjSubvariety.nativeProductPolynomialInclusion_coefficientMap}
\leanok
\uses{local:actual-original-product-scheme-projection}
The actual original product inclusion into polynomial projective space over
$k[t]$, followed by the actual projective coefficient map to $k$, equals the
original variety projection followed by the SAME original closed inclusion.
\end{lemma}
\begin{proof}
The genuine polynomial projective/base-product isomorphism has original
projective projection equal to the actual coefficient map. Compose its
inverse with that identity to recover the actual base-product projection.
The original product inclusion followed by that projection equals the
original variety projection followed by its actual closed inclusion.
Associativity proves the whole original scheme-morphism square.
\end{proof}

\begin{lemma}[The same normalized model retains the original polarization square]
\label{local:actual-normalized-polynomial-original-inclusion-square}
\lean{EnvelopeProperty.ProjSubvariety.normalizedSimultaneousPolynomialProjection_coefficientMap}
\leanok
\uses{local:actual-product-polynomial-original-inclusion-square,local:actual-normalized-original-projection}
The actual polynomial projective projection of the same normalized model,
followed by the actual coefficient map, equals its already constructed
original variety projection followed by the original closed inclusion.
\end{lemma}
\begin{proof}
The accepted original variety projection is literally normalization, original
simultaneous blowup, then original product projection. Substitute this proved
whole equality. The polynomial projection is normalization and original
blowup followed by the actual product polynomial inclusion. Compose the
preceding genuine product square with those original morphisms.
\end{proof}

\begin{theorem}[The exact original twist on the same normalized model]
\label{local:actual-normalized-polynomial-original-twist-iso}
\lean{EnvelopeProperty.ProjSubvariety.normalizedSimultaneousPolynomialTwistIso}
\leanok
\uses{projective:actual-original-pole-power-coefficient-pullback,local:actual-normalized-polynomial-original-inclusion-square}
On the same actual original normalized model, the actual pullback of the
original polynomial-base $\mathcal O(d)$ is isomorphic to the original
variety's SAME $\mathcal O(d)$ pulled back along the SAME original variety
projection. The original nonnegative exponent $d$ is retained exactly.
\end{theorem}
\begin{proof}
The whole coefficient-change theorem identifies the original $\mathcal O(d)$
over $k[t]$ with the genuine coefficient pullback of the original
$\mathcal O(d)$ over $k$. Pull its inverse comparison to the actual normalized
model. The genuine module pullback-composition isomorphism gives the pullback
along the polynomial projection followed by the coefficient map. Substitute
the preceding whole scheme square, then use the inverse pullback-composition
isomorphism for original variety projection and original closed inclusion.
The resulting inner pullback is precisely the accepted original twist on $X$.
\end{proof}

\begin{theorem}[The same actual polarized Cartier ratio line in the original class]
\label{local:actual-normalized-polynomial-original-ratio-line-iso}
\lean{EnvelopeProperty.ProjSubvariety.normalizedSimultaneousPolynomialRatioLineIso}
\leanok
\uses{local:actual-normalized-polynomial-original-twist-iso}
The actual polynomial-polarized ratio line whose genuine section weights were
proved convex is isomorphic to the accepted whole original normalized-model
line $p^*\mathcal O_X(d)\otimes L_{I/J}$ on the SAME model. The actual ratio
line and original degree are both retained; no polarization comparison is data.
\end{theorem}
\begin{proof}
Tensor the preceding actual whole original twist isomorphism with the identity
of the SAME actual canonical Cartier ratio line. The explicit actual tensor
comparison gives exactly the two originally defined polarized ratio lines.
\end{proof}
